\documentclass[11pt]{article}

\usepackage[margin=1in]{geometry}
\usepackage[T1]{fontenc}
\usepackage{lmodern}
\usepackage{microtype}
\usepackage{amsmath,amssymb,amsthm,mathtools}
\usepackage{booktabs,tabularx,array}
\usepackage[inline]{enumitem}
\usepackage[dvipsnames]{xcolor}
\usepackage[colorlinks=true,linkcolor=MidnightBlue,citecolor=BrickRed,
  urlcolor=MidnightBlue]{hyperref}

\newtheorem{theorem}{Theorem}[section]
\newtheorem{lemma}[theorem]{Lemma}
\newtheorem{proposition}[theorem]{Proposition}
\newtheorem{corollary}[theorem]{Corollary}
\newtheorem{conjecture}[theorem]{Conjecture}
\newtheorem{question}[theorem]{Question}
\theoremstyle{definition}
\newtheorem{definition}[theorem]{Definition}
\newtheorem{remark}[theorem]{Remark}

\newcommand{\bits}{\{0,1\}}
\newcommand{\F}{\mathbb F}
\newcommand{\C}{\mathrm{C}}
\newcommand{\poly}{\operatorname{poly}}
\newcommand{\MP}{\operatorname{MP}}
\newcommand{\GL}{\operatorname{GL}}
\newcommand{\proved}{\textcolor{ForestGreen}{\textbf{Proved}}}
\newcommand{\conditional}{\textcolor{MidnightBlue}{\textbf{Conditional}}}
\newcommand{\candidate}{\textcolor{BurntOrange}{\textbf{Candidate}}}
\newcommand{\open}{\textcolor{BrickRed}{\textbf{Open}}}

\title{\textbf{Polynomial-Size Mass Production of Boolean Circuits}\\[0.35em]
\large An exploratory transcript}
\author{Samuel Schlesinger}
\date{August 2026}

\begin{document}
\maketitle

\begin{abstract}
For a Boolean function $f$, let $f^{\times t}$ evaluate $f$ on $t$
independent input blocks.  This transcript asks for upper bounds on
$\C(f^{\times t})$ that use the actual, polynomial circuit size $\C(f)$ and
beat the replication bound $t\C(f)$.  There are three conclusions.

First, a theorem promising a strict saving for \emph{every} polynomial-size
function is false: functions of near-linear size are already forced to pay
linear input cost on every copy.  If $v(f)$ is the number of essential input
variables, the right optimistic target is instead
\[
  \C(f^{\times t})=O\bigl(\C(f)+t v(f)\bigr).
\]
This target is optimal up to a constant against elementary lower bounds, but
no unconditional proof is supplied here and none was found in the sources
audited below.

Second, Lipton proves the slightly weaker
$O(\C(f)\log\C(f)+tn)$ bound under the strong hypothesis that every
polynomial-time multiple-output Boolean map has linear-size circuits.
Applying his evaluator argument to the essential core replaces $n$ by $v(f)$.

Third, there is an unconditional polynomial-size example with a genuine
asymptotic saving.  For infinitely many $m$, a scalar bilinear function on
$2m$ essential bits has one-copy complexity $\Theta(m^2/\log m)$, whereas
$m$ independent copies have complexity $O(m^{\log_2 7})$.  Thus the batch
cost is $o(m\C(f))$.  The hard bilinear forms are obtained by counting, so
this is an existence result rather than an explicit family or a theorem for
all small circuits.
\end{abstract}

\section{Verdict and model}

\subsection{What has actually been obtained}

The exploration produced one unconditional positive theorem and a hierarchy
of universal conjectures.  The positive theorem is deliberately stated first
so that ``polynomial-size mass production'' does not refer only to a wish.

\begin{theorem}[Scalar polynomial-size mass production]
\label{thm:scalar-bilinear}
For every sufficiently large power of two $m$, there is an invertible matrix
$A_m\in\F_2^{m\times m}$ for which the scalar Boolean function
\[
  b_{A_m}(x,y)=y^{\mathsf T}A_mx,
  \qquad x,y\in\F_2^m,
\]
has all $2m$ variables essential and satisfies
\[
  \C(b_{A_m})=\Theta\!\left(\frac{m^2}{\log m}\right),
  \qquad
  \C(b_{A_m}^{\times m})=O\!\left(m^{\log_2 7}\right).
\]
Consequently,
\[
  \frac{\C(b_{A_m}^{\times m})}{m\C(b_{A_m})}
  =O\!\left(\frac{\log m}{m^{3-\log_2 7}}\right)=o(1).
\]
\end{theorem}

The proof is in Section~\ref{sec:bilinear}.  This is a single-output Boolean
function, not a vector-valued evasion of the question.  Its input length is
$2m$, its circuit size is quadratic up to a logarithm, and the number of
copies is linear, so all parameters are polynomial.  On the other hand, the
matrix $A_m$ is chosen nonconstructively.  This is a scalarization of the
standard matrix-multiplication mass-production example; no novelty or
priority claim is intended.

\paragraph{Status: \proved.}
This note is standalone for the polynomial-size analysis.  Its two references
to \texttt{main.tex} mean the affine-line construction and the
exponential-range mass-production theorem developed in that companion
manuscript.

\subsection{Circuit convention}

Circuits use a fixed finite complete Boolean basis of fan-in at most two,
such as $\{\wedge,\vee,\neg\}$.  Size counts non-input gates.  Constants,
fan-out, wires, and output designations are free.  XOR therefore has constant
cost.  Changing to any other fixed finite complete basis changes every bound
below by at most a constant factor.
The notation $\widetilde O$ suppresses factors polylogarithmic in the
parameters displayed in the same statement.

For $f:\bits^n\to\bits$ and $t\ge1$, define the $t$-fold independent product
\[
 f^{\times t}(x^{(1)},\ldots,x^{(t)})
   =\bigl(f(x^{(1)}),\ldots,f(x^{(t)})\bigr).
\]
The trivial upper bound is
\begin{equation}
  \C(f^{\times t})\le t\C(f).
  \label{eq:replication}
\end{equation}
The question is when the ratio between the two sides tends to zero.

\paragraph{Status vocabulary.}
\proved{} means that a proof is included or a cited theorem was checked.
\conditional{} means that the displayed conclusion depends on a stated
unproved hypothesis.  \candidate{} marks a proposed mechanism with a named
proof obligation.  \open{} means that no proof is claimed here.

\section{The intrinsic scale of the problem}
\label{sec:intrinsic}

\subsection{Essential arity}

\begin{definition}[Essential arity]
An input coordinate of $f:\bits^n\to\bits$ is \emph{essential} if changing
that coordinate alone can change $f$.  Write $v(f)$ for the number of
essential coordinates.  Restricting the inessential coordinates and keeping
one copy of every essential coordinate gives an \emph{essential core}
$f^{\rm core}:\bits^{v(f)}\to\bits$.
\end{definition}

\begin{lemma}[Dummy-padding invariance]
\label{lem:dummy}
For every $f$ and $t\ge1$,
\[
  \C(f)=\C(f^{\rm core}),
  \qquad
  \C(f^{\times t})=\C((f^{\rm core})^{\times t}).
\]
\end{lemma}

\begin{proof}
A circuit for the core computes $f$ after its input wires are relabelled and
the dummy pins are ignored.  Conversely, fixing all dummy inputs in a circuit
for $f$ produces a circuit of no larger size for the core.  Apply the same
two operations independently to every input block.
\end{proof}

Thus the ambient arity $n$ can be inflated without changing the task.  Every
size-sensitive statement below uses $v(f)$.

\begin{definition}[Individual mass-production curve]
\[
  \MP_f(t)=\C(f^{\times t}).
\]
\end{definition}

The elementary identities
\[
  \MP_f(1)=\C(f),\qquad
  \MP_f(t+u)\le\MP_f(t)+\MP_f(u),\qquad
  \MP_f(t)\le\MP_f(t+1)
\]
follow respectively from the definition, parallel composition, and fixing
one input block and discarding its output.  By subadditivity, Fekete's lemma
also gives the well-defined asymptotic cost per copy
\[
  \overline{\C}(f)
   :=\lim_{t\to\infty}\frac{\C(f^{\times t})}{t}
   =\inf_{t\ge1}\frac{\C(f^{\times t})}{t}.
\]

\subsection{A sharp elementary lower envelope}

The usual input-incidence argument loses a small amount when outputs are
allowed to be input wires.  The following graph form handles that case.

\begin{lemma}[Input--output incidence]
\label{lem:incidence}
Let $F:\bits^N\to\bits^q$ be a vector-valued Boolean function, and suppose
that $V$ input coordinates are essential for at least one output.  Then
\[
  \C(F)\ge V-q.
\]
\end{lemma}

\begin{proof}
Append one cost-free sink to every designated output of a circuit for $F$.
For each essential input, retain one directed path from that input to an
output it influences, and let $H$ be the union of these paths.  If $g_H$ is
the number of gates retained in $H$, $r\le q$ the number of retained output
sinks, $e$ the number of edges, and $c$ the number of weak components of $H$,
then every component contains a sink, so $c\le r$.  The underlying undirected
graph gives
\[
 e\ge V+g_H+r-c.
\]
At most $2g_H$ edges enter retained gates and exactly $r$ enter sinks, so
$e\le2g_H+r$.  Hence $g_H\ge V-c\ge V-r\ge V-q$.  The original circuit has
at least $g_H$ gates.
\end{proof}

\begin{proposition}[The unavoidable scale]
\label{prop:envelope}
For every scalar Boolean function $f$, with $s=\C(f)$ and $v=v(f)$, and
every $t\ge1$,
\begin{equation}
  \boxed{\max\{s,t(v-1)\}\le \C(f^{\times t})\le ts.}
  \label{eq:envelope}
\end{equation}
If $v\ge2$, then in particular
\[
  \C(f^{\times t})\ge\frac{s+tv}{3}.
\]
\end{proposition}

\begin{proof}
Fixing $t-1$ blocks and retaining one output yields a circuit for $f$, giving
the lower bound $s$.  The vector function $f^{\times t}$ has $tv$ essential
inputs and $t$ outputs, so Lemma~\ref{lem:incidence} gives $t(v-1)$.  The
upper bound is replication.  Finally, $v-1\ge v/2$ for $v\ge2$, whence
$s+tv\le3\max\{s,t(v-1)\}$.
\end{proof}

At $t=1$, the proposition also gives
\begin{equation}
  v(f)\le \C(f)+1.
  \label{eq:v-vs-s}
\end{equation}

\subsection{Why a uniform strict saving is impossible}

\begin{proposition}[Exact and constant-factor no-go examples]
\label{prop:nogo}
For $f(x_1,x_2)=x_1\wedge x_2$ over a basis containing AND,
\[
  \C(f)=1,\qquad \C(f^{\times t})=t.
\]
For parity on $v$ inputs over $\{\wedge,\vee,\neg\}$,
\[
  t(v-1)\le \C(\mathsf{PAR}_v^{\times t})\le4t(v-1).
\]
More generally, if every variable of $f_v$ is essential and
$\C(f_v)=O(v)$, then $\C(f_v^{\times t})=\Omega(t\C(f_v))$ uniformly in $t$.
\end{proposition}

\begin{proof}
For AND, the upper bound uses one gate per copy and the lower bound is
$t(v-1)=t$.  An XOR of two bits has a constant-size implementation, so a
chain computes parity in at most $4(v-1)$ gates; the lower bound again follows
from Proposition~\ref{prop:envelope}.  If $\C(f_v)\le Kv$, then
$t(v-1)\ge tv/2\ge t\C(f_v)/(2K)$ for $v\ge2$.
\end{proof}

Therefore a claim of $o(t\C(f))$ for \emph{all} polynomial-size functions is
literally false.  Along any asymptotic family with $v(f)\to\infty$ and
$\C(f)>0$, a strict saving can occur only if
\begin{equation}
  t\longrightarrow\infty
  \qquad\text{and}\qquad
  \frac{\C(f)}{v(f)-1}\longrightarrow\infty,
  \label{eq:necessary}
\end{equation}
equivalently $\C(f)/v(f)\to\infty$ in this regime.

\subsection{The right two targets}

\begin{conjecture}[Pointwise optimal mass production]
\label{conj:optimal}
There is a basis-dependent constant $K$ such that, for every scalar Boolean
function $f$ and every $t\ge1$,
\begin{equation}
  \C(f^{\times t})\le K\bigl(\C(f)+t v(f)\bigr).
  \label{eq:optimal-target}
\end{equation}
\end{conjecture}

\paragraph{Status: \open.}
Proposition~\ref{prop:envelope} shows that the right side is optimal up to a
factor of three.  It would give
\[
 \frac{\C(f^{\times t})}{t\C(f)}
   =O\!\left(\frac1t+\frac{v(f)}{\C(f)}\right),
\]
so the necessary conditions in \eqref{eq:necessary} would also be sufficient.

The first milestone allows the natural cost of describing a circuit.

\begin{conjecture}[Description-loss mass production]
\label{conj:description}
For every scalar Boolean function $f$ and every $t\ge1$,
\begin{equation}
 \C(f^{\times t})
 =O\!\left(\min\left\{t\C(f),\,
 \C(f)\log(\C(f)+2)+t v(f)\right\}\right).
 \label{eq:description-target}
\end{equation}
\end{conjecture}

\paragraph{Status: \open{} unconditionally.}
Writing $s=\C(f)$, this is a strict saving when
$t/\log(s+2)\to\infty$ and $s/v(f)\to\infty$.  If
$s=\Theta(v^a)$ and $t=\Theta(v^b)$ with $a\ge1$ and $b\ge0$, the optimal
target has size
\[
 O(v^a+v^{b+1})
\]
and improves on $v^{a+b}$ by a factor
$\Omega(v^{\min\{b,a-1\}})$.  The description-loss target gives the same
polynomial phase diagram, with a logarithmic loss.  Thus both predict a
strict saving exactly when $a>1$ and $b>0$.

\section{Three different meanings of a universal theorem}
\label{sec:strengths}

Fix a topologically ordered gate-list encoding.  A circuit with $v$ inputs
and $s$ gates has a description of length
\begin{equation}
 d(v,s)=O\bigl((s+1)\log(v+s+2)\bigr).
 \label{eq:description-length}
\end{equation}
For a minimum scalar circuit with at least two essential inputs,
\eqref{eq:v-vs-s} reduces this to $O(s\log(s+2))$.

There are three logically distinct goals.

\begin{enumerate}[label=\textnormal{(\roman*)}]
\item \emph{Pointwise existence:} for every $f$, some small circuit computes
  $f^{\times t}$, with the bound expressed using the minimum $\C(f)$.
\item \emph{Effective compilation:} an algorithm receives a particular
  circuit $Q$ of size $s$ and $t$ in unary, and outputs a small circuit for
  $(f_Q)^{\times t}$.  The bound is now in terms of the supplied $s$, not the
  unknown minimum $\C(f_Q)$.
\item \emph{Universal batch evaluation:} one circuit $U_{v,s,t}$ receives
  the description of any size-$s$ circuit together with $t$ data blocks and
  returns all $t$ evaluations.
\end{enumerate}

An effective universal evaluator gives an effective compiler by hardwiring
the program, and a compiler gives pointwise existence by applying it to a
minimum circuit as a logical witness.  A merely nonuniform universal circuit
need not be constructible, and separate compiled circuits need not assemble
into one small universal evaluator.  No reverse implication is automatic.

The $s\log s$ cost belongs naturally to the live program input of the third
formulation.  It is not an information-theoretic lower bound on the circuit
left after the program is hardwired.

\subsection{Lipton's conditional theorem}

Lipton's exact statement uses the declared arity $n$.  Lemma~\ref{lem:dummy}
lets us apply the same argument to the essential core and use $v(f)$ instead.
The following restricted form of Lipton's hypothesis is sufficient for that
evaluator argument; it is extremely strong and unproved.

\begin{quote}
\textbf{Hypothesis $\mathbf H_{\rm func}$.}
Every polynomial-time uniform family of Boolean maps
$G_N:\bits^N\to\bits^{m(N)}$, where $m(N)\le N$, has circuits of size
$O(N)$.  The hidden constant may depend on the family $G$.
\end{quote}

\begin{theorem}[Derived from Lipton's evaluator argument \cite{Lipton1994}]
\label{thm:lipton}
Assume $\mathbf H_{\rm func}$.  Then, for every scalar Boolean function $f$
and every $t\ge1$,
\[
  \C(f^{\times t})
  =O\bigl(t v(f)+\C(f)\log(\C(f)+2)\bigr).
\]
The same argument yields a nonuniform universal batch evaluator
of size $O(tv+d(v,s))$ for size-$s$ circuits.  It does not supply an
algorithm that constructs those evaluator circuits.
\end{theorem}

\begin{proof}
First handle $v\le1$ by replication.  For $v\ge2$, choose a standard
fixed-width gate-list field of exact padded length
$D(v,s)=\Theta((s+1)\log(v+s+2))$.  Define one canonical function
$E_N:\bits^N\to\bits^N$.  Its input contains a self-delimiting header
$(v,s,t,D)$, a $D$-bit circuit field, $t$ blocks of $v$ data bits, and padding
to total length
\[
  N=tv+D+O\bigl(\log(v+s+t+D+2)\bigr)=\Theta(tv+D),
\]
where the hidden constants come only from the fixed header format.  On a valid
description, $E_N$ simulates the encoded
circuit in topological order on every block, writes the $t$ answers in its
first $t$ output positions, and fills the rest with zeroes.  Invalid encodings
map to zero.  Here $s\le D\le N$ and $t\le tv\le N$, so naive $O(ts)$
simulation takes at most $O(N^2)$ time.  Thus the $E_N$ form one fixed
polynomial-time vector-function family.

Hypothesis $\mathbf H_{\rm func}$ gives $E_N$ circuits of size $O(N)$.
Hardwire the header and padding to obtain a nonuniform universal evaluator of
size $O(tv+d(v,s))$.  Finally hardwire a minimum circuit for $f$ and use
\eqref{eq:description-length} and \eqref{eq:v-vs-s}.  Constants and
one-variable functions were handled at the start.
\end{proof}

\paragraph{Status: \conditional.}
The evaluator algorithm is explicit; the small circuit implementing it is
conditional and nonconstructive.  Valiant's universal circuits
\cite{Valiant1976} explain the ordinary $O(s\log s)$ programmable-simulation
cost, but do not themselves remove the factor $t$ from $t$ evaluations.

\section{What the worst-case mass-production theorems do not yet give}
\label{sec:literature}

The checked classical results answer the question when the one-copy function
is allowed to have worst-case exponential complexity.  They do not give the
desired dependence on an actual polynomial value $s=\C(f)$.  This is a claim
about the sources audited below, not an exhaustive no-result claim.

\subsection{Checked baseline bounds}

Shannon's counting lower bound and Lupanov's synthesis upper bound give the
worst-case one-output scale
\begin{equation}
  \max_{f:\bits^v\to\bits}\C(f)=\Theta(2^v/v)
  \label{eq:shannon-lupanov}
\end{equation}
for any fixed complete finite basis; see
\cite{Shannon1949,Lupanov1958,Wegener1987}.  Against that scale, the following
are genuine mass-production results.

\begin{itemize}
\item Let
$L_t(v)=\max_{f:\bits^v\to\bits}\C(f^{\times t})$.  In a fixed complete
weighted basis with Lupanov constant $\rho$, Uhlig proved that every sequence
$t=t(v)$ satisfying $\log t=o(v/\log v)$ obeys
\[
  L_t(v)\sim\rho\frac{2^v}{v}.
\]
Consequently the $O(2^v/v)$ upper bound holds uniformly for every $f$, but
the matching lower bound is worst-case only and does not compare the batch
size with an individual function's $\C(f)$.  A modern restatement and proof
overview appear in \cite{Kretschmer2023}; see also
\cite{Uhlig1974,Uhlig1992,Wegener1987}.

\item Paul's Lemma~2 \cite{Paul1976} gives, in his arbitrary-binary-gate
convention,
\begin{equation}
  \C(f^{\times t})
    =O\!\left(\max\{v2^v,\,vt\log^2(t+2)\}\right).
  \label{eq:paul}
\end{equation}
A constant-factor basis simulation transfers the estimate to the present
model.

\item The companion construction in \texttt{main.tex} proves the
order-of-growth bound
\[
  \C(f^{\times t})=O_\gamma(2^v/v)
  \qquad(t\le2^{\gamma v})
\]
for every fixed $\gamma<1$.  Unlike Uhlig's smaller-range theorem, that
argument does not retain the optimal one-copy leading constant.
\end{itemize}

Paul's bound also explains the known polynomial amortized cost when the copy
count is allowed to become very large.  Partition requests into blocks of
$k_0=\lceil2^v/v^2\rceil$ and apply \eqref{eq:paul} to every block.  This
gives the following derived consequence, valid for every $f$ and $t$:
\begin{equation}
  \C(f^{\times t})=O\bigl(v^3t+v2^v\bigr).
  \label{eq:paul-blocked}
\end{equation}
Thus the asymptotic cost per copy is $O(v^3)$.  The setup term is still
exponential, so this only helps after an exponential number of requests.

Two other statements in Paul's paper must be kept separate.  His scalar
two-copy theorem gives a small circuit for the \emph{single-output} function
$f(x)\vee f(y)$; it is not a circuit for the two-output map $f\times f$.
His genuine two-output direct-product theorem is for vector-valued functions.
Paul explicitly records that the corresponding nontrivial scalar
$f\times f$ saving was not obtained there.

Albrecht's short paper \cite{Albrecht1989} studies simultaneous realization.
Its publisher preview and abstract promise small size and depth overhead for
the corresponding Shannon functions.  This transcript does not use a sharper formula from that
paper: a theorem-level primary-source comparison has not been recovered here.
This omission is intentional rather than an invitation to infer a bound.

\subsection{The polynomial-size gap}

Combining only the checked generic estimates gives the envelope
\begin{equation}
 \C(f^{\times t})\le
 \min\left\{
   t\C(f),\;
   O\!\left(\max\{v2^v,vt\log^2(t+2)\}\right),\;
   O\bigl(v^3t+v2^v\bigr)
 \right\},
 \label{eq:known-envelope}
\end{equation}
with the $O(2^v/v)$ plateau available in the copy ranges just described.
If both $\C(f)$ and $t$ are polynomial in $v$, every displayed generic
alternative to replication retains an exponential setup term.  Hence these
theorems alone return $t\C(f)$ as the useful bound in precisely the regime of
interest.  They do not refute a size-sensitive theorem; they simply do not
interpolate to the actual value of $\C(f)$.

\subsection{Why ordinary counting does not prove a direct sum}

The number of fan-in-two circuits of size at most $q$ with $N$ inputs and
one output is
\begin{equation}
  2^{O((q+1)\log(N+q+2))}.
  \label{eq:circuit-count}
\end{equation}
Indeed, for every topologically ordered gate one chooses a constant-size gate
type and at most two predecessors among $N+q+2$ available sources, followed
by one output source.  With $r$ designated outputs the corresponding upper
count is
\begin{equation}
  2^{O((q+r+1)\log(N+q+2))},
  \label{eq:multioutput-count}
\end{equation}
because one must also list the $r$ output sources.

The map $f\mapsto f^{\times t}$ is injective, but it introduces no new choice:
\[
  \left|\{f^{\times t}:f:\bits^v\to\bits\}\right|=2^{2^v},
\]
the same cardinality as the one-copy family.  Likewise, the products of
functions having circuits of size below $2s$ form a set of size at most
$2^{O((s+1)\log(v+s+2))}$, with no factor $t$ in the exponent.  Applying
\eqref{eq:multioutput-count} to candidate batch circuits adds an output-list
term $t\log(tv+q+2)$ to the \emph{circuit} count but adds no entropy to the
target family; this only weakens a counting lower bound.  In parameter ranges
where output listing is not already dominant, plain Shannon counting can
force an $s$-scale lower bound, but it cannot manufacture a $ts$-scale direct
sum merely by replacing every target with its $t$-fold product.  A stronger
lower bound for a particular $f$ would need structural information, not just
target-set entropy.

\section{An unconditional scalar pilot theorem}
\label{sec:bilinear}

This section proves Theorem~\ref{thm:scalar-bilinear}.  The construction is
the cleanest evidence found here that polynomial-size scalar functions can
exhibit the desired phenomenon.

\subsection{The standard vector warmup}

For a fixed matrix $A\in\F_2^{m\times m}$, the linear map $x\mapsto Ax$ has
complexity $\Omega(m^2/\log m)$ for almost every $A$ by counting.  Evaluating
it on $m$ vectors at once is the matrix product $AX$, where the requested
vectors are the columns of $X$.  Fast matrix multiplication computes this in
$O(m^\omega)$ Boolean gates for some $\omega<3$; this standard example is
recalled in \cite{Kretschmer2023}.  The only defect for the present question
is that $x\mapsto Ax$ has $m$ output bits.  A bilinear outer shell turns it
into a scalar function without destroying the saving.

\subsection{One copy is genuinely quadratic}

For $A\in\F_2^{m\times m}$, put
\[
  b_A(x,y)=y^{\mathsf T}Ax
    =\bigoplus_{i,j\in[m]} A_{ij}y_i x_j.
\]

\begin{lemma}[Almost every invertible form is hard]
\label{lem:bilinear-hard}
There is a constant $c>0$ such that, for all sufficiently large $m$, all but
an exponentially small fraction of the matrices in $\GL(m,2)$ satisfy
\[
  \C(b_A)\ge c\frac{m^2}{\log m}.
\]
Every such $b_A$ has all $2m$ variables essential.
\end{lemma}

\begin{proof}
The map $A\mapsto b_A$ is injective, because
$b_A(e_j,e_i)=A_{ij}$.  Moreover,
\[
 |\GL(m,2)|
   =2^{m^2}\prod_{j=1}^m(1-2^{-j})
   \ge c_0 2^{m^2}
\]
for an absolute constant $c_0>0$.  By \eqref{eq:circuit-count}, the number of
functions on $2m$ inputs computed with at most
$q=c m^2/\log m$ gates is $2^{O(c m^2)}$.  Choosing $c$ sufficiently small
makes this $2^{m^2-\Omega(m^2)}$, hence a vanishing fraction of the invertible
forms.

If $A$ is invertible, every row and every column is nonzero.  A nonzero entry
$A_{ij}$ witnesses that $x_j$ is essential after $y=e_i$ and the other
$x$-coordinates are fixed to zero; symmetrically it witnesses that $y_i$ is
essential after $x=e_j$.
\end{proof}

\begin{lemma}[Four-Russians upper bound]
\label{lem:bilinear-upper}
For every $A\in\F_2^{m\times m}$,
\[
  \C(b_A)=O(m^2/\log m).
\]
\end{lemma}

\begin{proof}
Let $k=\lfloor\log_2m\rfloor$ and partition the coordinates of $x$ into
$O(m/k)$ blocks of length at most $k$.  For one block, all $2^k$ linear forms
in its bits can be generated using $O(2^k)$ XOR gates: recursively retain all
forms on the first $k-1$ bits and XOR each of them with the last bit.  Across
all blocks this costs
\[
  O\bigl((m/k)2^k\bigr)=O(m^2/\log m).
\]
For every row of $A$, select the required precomputed form from each block and
XOR the $O(m/k)$ selected values.  All $m$ rows of $Ax$ cost another
$O(m^2/k)$.  Finally compute the $m$ products $y_i(Ax)_i$ and XOR them, at
cost $O(m)$.  Replacing each XOR with a constant number of basis gates proves
the claim.
\end{proof}

Together, Lemmas~\ref{lem:bilinear-hard} and~\ref{lem:bilinear-upper} show
that almost every invertible $A$ has
$\C(b_A)=\Theta(m^2/\log m)$.

\subsection{Many copies become one matrix product}

\begin{proposition}[Batching bilinear forms]
\label{prop:bilinear-batch}
For every fixed $A\in\F_2^{m\times m}$ and every $t\ge1$,
\[
 \C(b_A^{\times t})
   =O\!\left(\left\lceil\frac tm\right\rceil
       m^{\log_2 7}+tm\right).
 \label{eq:bilinear-batch}
\]
\end{proposition}

\begin{proof}
First take $t=m$.  Arrange the $x$-inputs as the columns of an $m\times m$
matrix $X$.  Strassen's seven-product recursion \cite{Strassen1969} computes
$Z=AX$ over $\F_2$ using $O(m^{\log_2 7})$ field operations.  It remains to
compute $y_r^{\mathsf T}Z_{*,r}$ for every column $r$, which takes $O(m)$
AND/XOR gates per column and $O(m^2)$ in total.  Since $\log_2 7>2$, the
matrix multiplication dominates.  Hardwiring $A$ into a circuit for general
matrix multiplication cannot increase its size.

For arbitrary $t$, divide the requests into groups of at most $m$, pad the
last group with zero columns, and apply the same circuit to every group.  The
inner products over all genuine requests cost $O(tm)$.
\end{proof}

Choosing a hard invertible $A_m$ from Lemma~\ref{lem:bilinear-hard} and
setting $t=m$ proves Theorem~\ref{thm:scalar-bilinear}.  More generally,
\eqref{eq:bilinear-batch} yields a strict saving for every matrix family for
which
\[
  \C(b_{A_m})=\omega(m^{\log_2 7-1})
\]
and every sequence $t=t(m)\ge m$.  Indeed,
$\lceil t/m\rceil\le2t/m$, so the ratio in question is
\[
 O\!\left(
  \frac{m^{\log_2 7-1}+m}{\C(b_{A_m})}
 \right)=o(1).
\]
Counting supplies many such matrices, but proving that bound for an explicit
$A_m$ would itself prove an unrestricted superlinear circuit lower bound.  No
such lower bound is proved here.  The nonconstructivity is therefore a real
boundary, not a cosmetic omission.

\section{Mechanism audit for a theorem about every small circuit}
\label{sec:mechanisms}

The bilinear theorem identifies a successful mechanism: a large amount of
one-copy work is a fixed linear transformation, and different instances can
be mixed and later separated by a fast algebraic algorithm.  An all-functions
theorem would need either to find an analogue of that structure in every
superlinear circuit or to use a different kind of global reorganization.

\subsection{Gate-by-gate SIMD is not mass production}

Given a size-$s$ circuit $Q$, one can share its topology while maintaining a
length-$t$ vector of values at every gate.  At the bit-gate level this still
uses $\Theta(t)$ gates for each genuinely binary operation, hence
$\Theta(ts)$ gates overall.  Proposition~\ref{prop:nogo} shows that even
$\mathsf{AND}_2^{\times t}$ costs exactly $t$.  Sharing a wiring diagram is
not the same as sharing the values computed on independent inputs.

\subsection{Macro tabulation needs small interfaces}

A tempting Four-Russians plan partitions $Q$ into $k$-gate macros, tabulates
each macro, and answers all requests by table lookup.  Let $b$ be the number
of live input wires crossing into a macro and $r$ the number of values leaving
it.  An arbitrary $b$-input, $r$-output table already has description scale
$r2^b$, and the $t$ requests contain $tb$ independent live input bits.  If
$b=\Theta(k)$, tabulation is exponentially more expensive than the macro it
replaces.  Choosing $k\asymp\log t$ does not by itself fix the problem: the
boundary cost and the per-query routing must both be charged.

There is a precise restricted criterion behind this idea.  Partition the
gates of $Q$ into modules $M_i$ compatible with the quotient DAG, and let
$\phi_i:\bits^{b_i}\to\bits^{r_i}$ be the function taking the distinct live
values entering $M_i$ to the distinct values leaving it.  Module substitution
gives
\begin{equation}
  \C(f^{\times t})\le\sum_i\C(\phi_i^{\times t}).
  \label{eq:module-substitution}
\end{equation}
Indeed, replace every module, in quotient-topological order, by a circuit for
$t$ independent copies of its interface map, keeping instance indices aligned
on every connecting wire.  This observation is exact, but replication on the
right side gives
$\C(\phi_i^{\times t})\le t\C(\phi_i)\le t|M_i|$ and hence recovers the
upper bound $ts$.  Decomposition only helps if a
substantial collection of the interface maps is itself batchable.

\paragraph{Load-bearing lemma.}
A decomposition proof would need a statement that every \emph{minimum}
superlinear circuit has enough large pieces with boundary much smaller than
their internal size, together with a batch-lookup circuit whose total routing
cost is below $tk$ per piece.  Bounded-degree layered expanders, oriented and
reduced to fan-in two by constant-size gadgets, are a natural adversarial test:
any proposed lemma must explain why their medium-size pieces do not expose
linear boundary.  No general decomposition lemma is proved here.

\subsection{Restriction coding is not automatically size-sensitive}

The construction in \texttt{main.tex} reorganizes the truth table of an
arbitrary function by coding many prefix restrictions and recovering requested
ones along affine lines.  Its resource functions are restrictions and linear
combinations of restrictions on a shorter suffix.  Worst-case synthesis
charges them by their suffix arity and gives the desired $2^v/v$ scale.

Starting instead from a size-$s$ circuit does not give a corresponding bound
on the \emph{sum} of the minimum sizes of all those derived resource
functions.  Synthesizing every restriction separately can cost $2^p s$ after
fixing a $p$-bit prefix, before the scheduling recursion even begins.  The
missing statement is therefore not another choice of code parameters; it is
a restriction-sharing theorem that relates the aggregate complexity of a
decodable resource family to the original $s$.

Even the simplest path-difference charge fails exponentially.

\begin{proposition}[Restriction differences need not charge to the circuit]
\label{prop:restriction-obstruction}
For every sufficiently large $k$, there are $m=32k$ and a linear map
$G:\F_2^k\to\F_2^m$ such that the polynomial-size scalar function
\[
  F_G(a,z)=\langle Ga,z\rangle,
  \qquad a\in\F_2^k,\quad z\in\F_2^m,
  \qquad \C(F_G)=O(km),
\]
has the following property.  Write
$(F_G)_a(z):=F_G(a,z)$ for a prefix restriction.  For every ordering $\pi$
of its $2^k$ prefix restrictions,
\[
 \sum_{i=1}^{2^k-1}
 \C\bigl((F_G)_{\pi(i-1)}\oplus(F_G)_{\pi(i)}\bigr)
 \ge(2^k-1)(m/4-1).
\]
\end{proposition}

\begin{proof}
A random $m\times k$ binary matrix has, with positive probability,
$\mathrm{wt}(Gu)\ge m/4$ for every nonzero $u$: for fixed $u$, $Gu$ is
uniform, a Chernoff bound gives failure probability at most $e^{-m/16}$, and
a union bound over fewer than $2^k$ choices succeeds when $m=32k$.  Compute
$Ga$ and then its inner product with $z$ to get the $O(km)$ upper bound.

The difference between restrictions at distinct $a,a'$ is the parity with
coefficient vector $G(a-a')$.  It depends on at least $m/4$ variables, so
\eqref{eq:v-vs-s} gives circuit size at least $m/4-1$.  Every consecutive pair
in an ordering is distinct, proving the sum.
\end{proof}

This refutes a universal assertion that consecutive restriction differences
have total complexity $\poly(k)\C(f)$.  On the other hand, the $k$ functions
$h_j(z)=\langle Ge_j,z\rangle$ generate the restrictions by
$(F_G)_a=\bigoplus_{j:a_j=1}h_j$ and have aggregate size $O(km)$.  If the
$j$th address bit is one in all $t$ requests, however, $h_j$ is demanded on
all $t$ unrelated suffixes.  Thus a viable code must control resource
complexity and request congestion simultaneously.

\paragraph{Status: \candidate.}
Here is a formal template for the missing lemma.  For a split input
$x=(a,z)\in\bits^p\times\bits^{v-p}$, an XOR restriction scheme consists of
resource functions $\mathcal H=\{h_1,\ldots,h_M\}$ and, for every address
$a$, a nonempty collection $\mathcal R(a)$ of recovery sets such that
\[
 f(a,z)=\bigoplus_{j\in R}h_j(z)
 \qquad\text{for every }R\in\mathcal R(a).
\]
For $t$ runtime addresses, a selector chooses $R_i\in\mathcal R(a_i)$; its
load on resource $j$ is $d_j=|\{i:j\in R_i\}|$.  The selector must also route
the corresponding suffixes into fixed-capacity resource circuits and route
their answers back to the $t$ decoders.

\begin{question}[Candidate restriction-resource budget]
\label{q:restriction-budget}
For some growing $1\le p=p(s,t)\le v$, does every size-$s$ function admit an
XOR restriction scheme, load caps $D_1,\ldots,D_M$, and a selector/router of
size $\widetilde O(tv)$ such that, for every tuple of $t$ runtime addresses,
it chooses recoveries with $d_j\le D_j$ and
\[
  \sum_{j=1}^M\C(h_j^{\times D_j})
    =\widetilde O(s+tv)?
\]
\end{question}

The quantifiers ``every $f$'' and ``aggregate'' are essential.  A proposal
that bounds each $\C(h)$ separately by $s$ merely recreates the factor being
removed.  Multiplexer circuits with many unrelated branches and circuits
whose restricted gate functions rapidly diversify are immediate stress
tests for Question~\ref{q:restriction-budget}.

\subsection{Circuit-as-static-data compilation}

The most conceptually direct route is to treat the circuit description as a
one-time database and the $t$ input blocks as queries.  This is exactly the
universal evaluator in Section~\ref{sec:strengths}.  Lipton's argument proves
that a linear-circuit hypothesis for polynomial-time maps would compress the
whole computation to
\[
  O(s\log s+tv).
\]
Unconditionally, ordinary universal circuits charge roughly $s\log s$ for
\emph{each separately simulated evaluation} unless some additional batch
structure is found.  The P-completeness of the circuit value problem
\cite{Ladner1975} is a warning about uniform parallel evaluation, but it is
not a nonuniform circuit lower bound and therefore is not evidence against
Conjectures~\ref{conj:optimal} or~\ref{conj:description}.

\subsection{Algebraic cores and other restricted classes}

Theorem~\ref{thm:scalar-bilinear} suggests looking for circuits with an
expensive static algebraic core: linear maps, bilinear maps, or low-degree
tensor contractions followed by only $O(v)$ instance-specific work.  Fast
rectangular matrix multiplication could extend the copy range in such
classes.  What is missing for general Boolean circuits is a normal-form
theorem that exposes a sufficiently large algebraic core without increasing
the circuit size by the savings one hopes to obtain.

Other plausible pilot restrictions are small separator number, bounded
congestion decompositions, or an explicit promise that every large module has
polylogarithmic interface.  Results for those classes would be informative,
but they should be stated in terms of the promised structural parameter; they
would not by themselves establish a theorem for arbitrary size-$s$ DAGs.

\section{The next theorem to try}
\label{sec:next}

The clean all-parameter conjectures may be stronger than necessary for a
first result.  The following polynomial-band formulation isolates a clean
polynomially separated regime in which a uniform asymptotic saving is
possible.

\begin{conjecture}[Polynomial-band milestone]
\label{conj:polynomial-band}
For every fixed $A>1$, $B>0$, and
$0<\varepsilon<\min\{A-1,B\}$, there is a constant
$K=K(A,B,\varepsilon)$ such that the following holds for all sufficiently
large $v$.  If every input of $f:\bits^v\to\bits$ is essential,
\[
 v^{1+\varepsilon}\le s=\C(f)\le v^A,
 \qquad
 v^\varepsilon\le t\le v^B,
\]
then
\[
 \C(f^{\times t})
 \le K\bigl(s\log(s+2)+tv\bigr).
\]
\end{conjecture}

This would imply uniformly over the promised band that
\[
 \frac{\C(f^{\times t})}{ts}
 =O\!\left(\frac{\log v}{v^\varepsilon}+\frac1{v^\varepsilon}\right)
 =o(1).
\]
It is weaker than Conjecture~\ref{conj:description}, avoids the regimes where
the elementary lower bounds forbid a saving, and directly answers the
polynomial-size question.  It is still \open.

The most concrete next candidate is Question~\ref{q:restriction-budget}.
It has three advantages over the raw conjecture: it interfaces directly with
the already developed recovery-code construction; it exposes every cost that
must be paid; and it has obvious counterexample searches.  A useful first
experiment would compute, for small optimized circuits and several prefix
sets, the number and aggregate circuit size of distinct restricted gate
functions.  A negative experiment would not refute all mass production, but
it would quickly kill the current route.

\section{Results ledger}

\begin{center}
\small
\renewcommand{\arraystretch}{1.18}
\begin{tabularx}{\linewidth}{@{}
  >{\raggedright\arraybackslash}p{0.19\linewidth}
  >{\raggedright\arraybackslash}p{0.20\linewidth}
  >{\raggedright\arraybackslash}X@{}}
\toprule
Status & Statement & Evidence or boundary \\
\midrule
\proved &
$\max\{\C(f),t(v(f)-1)\}\le\C(f^{\times t})\le t\C(f)$ &
Proposition~\ref{prop:envelope}; the lower bound is a fan-in graph count. \\

\proved &
No theorem can guarantee a strict saving for every polynomial-size function &
$\mathsf{AND}_2$ is exact and every essentially $v$-ary, $O(v)$-size family
has a constant-factor direct sum; Proposition~\ref{prop:nogo}. \\

\proved &
Scalar polynomial-size mass production exists &
Almost every invertible bilinear form has size $\Theta(m^2/\log m)$, and $m$
copies cost $O(m^{\log_2 7})$; Theorem~\ref{thm:scalar-bilinear}. \\

\proved &
Worst-case truth-table mass production &
Uhlig, Paul, and the companion theorem give strong exponential-size
plateaux, but their checked bounds retain exponential setup in the
polynomial-size regime. \\

\conditional &
$O(\C(f)\log\C(f)+tv(f))$ for every $f$ &
Theorem~\ref{thm:lipton}, assuming $\mathbf H_{\rm func}$.  This proves
nonuniform existence, not an effective compiler. \\

\open &
$O(\C(f)+tv(f))$ for every $f$ &
Conjecture~\ref{conj:optimal}; pointwise optimal up to a constant against the
elementary lower envelope. \\

\open &
An unconditional polynomial-band theorem &
Conjecture~\ref{conj:polynomial-band}; the proposed next route is the
restriction-resource budget in Question~\ref{q:restriction-budget}. \\
\bottomrule
\end{tabularx}
\end{center}

\paragraph{Bottom line.}
This transcript proves a polynomial-size mass-production theorem for a large
nonconstructive family of scalar Boolean functions, with a polynomial number
of copies and a vanishing ratio to replication.  The genuinely interesting
universal theorem is left open in this transcript, and no proof was found in
the audited sources.  Its plausible optimal form is
$O(\C(f)+tv(f))$; its first description-sensitive form is
$O(\C(f)\log\C(f)+tv(f))$.  Any proof has to reorganize computations across
independent inputs, not merely share circuit syntax.

\begin{thebibliography}{99}

\bibitem{Albrecht1989}
Andreas Albrecht.
\newblock On simultaneous realizations of Boolean functions, with applications.
\newblock In \emph{Parcella '88}, Lecture Notes in Computer Science 342,
pages 51--56. Springer, 1989.
\newblock \href{https://doi.org/10.1007/3-540-50647-0_102}{doi:10.1007/3-540-50647-0\_102}.

\bibitem{Kretschmer2023}
William Kretschmer.
\newblock Quantum mass production theorems.
\newblock In \emph{18th Conference on the Theory of Quantum Computation,
Communication and Cryptography}, LIPIcs 266, Article 10, 2023.
\newblock \href{https://doi.org/10.4230/LIPIcs.TQC.2023.10}{doi:10.4230/LIPIcs.TQC.2023.10}.

\bibitem{Ladner1975}
Richard E. Ladner.
\newblock The circuit value problem is log space complete for $P$.
\newblock \emph{SIGACT News}, 7(1):18--20, 1975.
\newblock \href{https://doi.org/10.1145/990518.990519}{doi:10.1145/990518.990519}.

\bibitem{Lipton1994}
Richard J. Lipton.
\newblock Some consequences of our failure to prove non-linear lower bounds
on explicit functions.
\newblock In \emph{Proceedings of the Ninth Annual IEEE Conference on
Structure in Complexity Theory}, pages 79--87, 1994.
\newblock \href{https://doi.org/10.1109/SCT.1994.315815}{doi:10.1109/SCT.1994.315815}.

\bibitem{Lupanov1958}
Oleg B. Lupanov.
\newblock On a method of circuit synthesis.
\newblock \emph{Izvestiya VUZ, Radiofizika}, 1(1):120--140, 1958.
\newblock In Russian.

\bibitem{Paul1976}
Wolfgang J. Paul.
\newblock Realizing Boolean functions on disjoint sets of variables.
\newblock \emph{Theoretical Computer Science}, 2(3):383--396, 1976.
\newblock \href{https://doi.org/10.1016/0304-3975(76)90089-X}{doi:10.1016/0304-3975(76)90089-X}.

\bibitem{Shannon1949}
Claude E. Shannon.
\newblock The synthesis of two-terminal switching circuits.
\newblock \emph{Bell System Technical Journal}, 28(1):59--98, 1949.
\newblock \href{https://doi.org/10.1002/j.1538-7305.1949.tb03624.x}{doi:10.1002/j.1538-7305.1949.tb03624.x}.

\bibitem{Strassen1969}
Volker Strassen.
\newblock Gaussian elimination is not optimal.
\newblock \emph{Numerische Mathematik}, 13(4):354--356, 1969.
\newblock \href{https://doi.org/10.1007/BF02165411}{doi:10.1007/BF02165411}.

\bibitem{Uhlig1992}
Dietmar Uhlig.
\newblock Networks computing Boolean functions for multiple input values.
\newblock In M. S. Paterson, editor, \emph{Boolean Function Complexity},
London Mathematical Society Lecture Note Series 169, pages 165--173.
Cambridge University Press, 1992.
\newblock \href{https://doi.org/10.1017/CBO9780511526633.013}{doi:10.1017/CBO9780511526633.013}.

\bibitem{Uhlig1974}
Dietmar Uhlig.
\newblock On the synthesis of self-correcting schemes from functional
elements with a small number of reliable elements.
\newblock \emph{Mathematical Notes of the Academy of Sciences of the USSR},
15(6):558--562, 1974.
\newblock \href{https://doi.org/10.1007/BF01152835}{doi:10.1007/BF01152835}.

\bibitem{Valiant1976}
Leslie G. Valiant.
\newblock Universal circuits (preliminary report).
\newblock In \emph{Proceedings of the Eighth Annual ACM Symposium on Theory
of Computing}, pages 196--203, 1976.
\newblock \href{https://doi.org/10.1145/800113.803649}{doi:10.1145/800113.803649}.

\bibitem{Wegener1987}
Ingo Wegener.
\newblock \emph{The Complexity of Boolean Functions}.
\newblock Wiley--Teubner, 1987.
\newblock \url{https://eccc.weizmann.ac.il/static/books/The_Complexity_of_Boolean_Functions/}.

\end{thebibliography}

\end{document}
